\rSec3[concept.booleantestable]{Boolean testability}

\begin{codeblock}
namespace std {
  template<class T>
  concept @\exposid{boolean-testable-impl}@ = convertible_to<T, bool>; // \expos
}
\end{codeblock}

\pnum
Let \tcode{e} be an expression such that \tcode{decltype((e))} is \tcode{T}. \tcode{T} models \exposconcept{boolean-testable-impl} only if
\begin{itemize}
\item either \tcode{remove_cvref_t<T>} is not a class type, or a search for the names \tcode{operator\&\&} and \tcode{operator||} in the scope of \tcode{remove_cvref_t<T>} finds nothing; and
\item argument-dependent lookup\iref{basic.lookup.argdep} for the names \tcode{operator\&\&} and \tcode{operator||} with \tcode{T} as the only argument type finds no disqualifying declaration (defined below).
\end{itemize}

\pnum
A \defnadj{disqualifying}{parameter} is a function parameter whose declared type \tcode{P}
\begin{itemize}
\item is not dependent on a template parameter, and there exists an implicit conversion sequence\iref{over.best.ics} from \tcode{e} to \tcode{P}; or
\item is dependent on one or more template parameters, and either
\begin{itemize}
\item \tcode{P} contains no template parameter that participates in template argument deduction\iref{temp.deduct.type}, or
\item template argument deduction using the rules for deducing template arguments in a function call\iref{temp.deduct.call} and \tcode{e} as the argument succeeds.
\end{itemize}
\end{itemize}

\pnum
A \defnadj{key}{parameter} of a function template \tcode{D} is a function parameter of type \cv{} \tcode{X} or reference thereto, where \tcode{X} names a specialization of a class template that has the same innermost enclosing non-inline namespace as \tcode{D}, and \tcode{X} contains at least one template parameter that participates in template argument deduction.
\begin{example}
In

\begin{codeblock}
namespace Z {
  template<class> struct C {};
  template<class T>
    void operator&&(C<T> x, T y);
  template<class T>
    void operator||(C<type_identity_t<T>> x, T y);
}
\end{codeblock}

the declaration of \tcode{Z::operator\&\&} contains one key parameter, \tcode{C<T> x}, and the declaration of \tcode{Z::operator||} contains no key parameters.
\end{example}

\pnum
A \defnadj{disqualifying}{declaration} is
\begin{itemize}
\item a (non-template) function declaration that contains at least one disqualifying parameter; or
\item a function template declaration that contains at least one disqualifying parameter, where
\begin{itemize}
\item at least one disqualifying parameter is a key parameter; or
\item the declaration contains no key parameters; or
\item the declaration declares a function template to which no name is bound\iref{dcl.meaning}.
\end{itemize}
\end{itemize}

\pnum
\begin{note}
The intention is to ensure that given two types \tcode{T1} and \tcode{T2} that each model \exposconcept{boolean-testable-impl}, the \tcode{\&\&} and \tcode{||} operators within the expressions \tcode{declval<T1>() \&\& declval<T2>()} and \tcode{declval<T1>() || declval<T2>()} resolve to the corresponding built-in operators.
\end{note}
